% Conclusiones (consigna 2.2, ítem 10). BORRADOR ASISTIDO POR IA (2026-11-12): reescribir con tus palabras.
\section{Conclusiones}

La oferta de financiamiento participativo financiero en el Perú creció con fuerza entre 2022 y 2026, de S/~\cifra{monto2022} a S/~\cifra{monto2025} millones anuales, pero está muy concentrada: solo \cifra{plataformasConFinanciamiento} de \cifra{plataformasAutorizadas} administradoras financiaron proyectos y una reúne el \cifra{participacionLider}~\% del monto. Es, sobre todo, un mercado de préstamos a empresas con montos crecientes, en el que solo la plataforma líder es rentable y retiene alrededor de \cifra{costoMediana}~\% de lo que financia. Las diferencias de colocación entre receptores reflejan el modelo de negocio de cada plataforma, y el mayor monto tras la modificación de 2025 es coherente con el crecimiento previo, por lo que no puede atribuirse a la norma.

Frente a un crédito tradicional de tasas altas, el FPF no puede compararse en precio: mientras no se publiquen la tasa y el plazo de sus operaciones, no podrá saberse si ofrece mejores condiciones que la banca y las cajas.
